% ===================================================================
\section{Tareas 1 a 4: de la red a las cascadas}

\diapo{7}{Tarea 1: red de pagos e índice de amplificación}{1.5 min}
\begin{decir}
La primera tarea construye la red y estudia solo dos términos: la autocorrección y el
contagio.

Generamos redes con tres formas. En la de Erdős–Rényi, cada par de entidades se conecta
al azar con la misma probabilidad; sirve de referencia. En la de Barabási–Albert, las
entidades nuevas prefieren conectarse con las que ya tienen muchas conexiones; aparecen
unos pocos nodos muy conectados. En la núcleo-periferia, un grupo pequeño de grandes
entidades está muy conectado entre sí y las demás se conectan con ese núcleo; es la
forma que se observa en los sistemas de pagos reales. Los montos siguen una distribución
lognormal: muchos pagos pequeños y pocos muy grandes.

Con solo estos dos términos, la ecuación es lineal: $\dot x=(-\alpha I+\beta W)x$. Como
todos los pagos son positivos, el teorema de Perron–Frobenius garantiza que el valor
propio dominante de $W$ es real y positivo; lo llamamos radio espectral $\rho(W)$. El
sistema es estable, es decir, el déficit se extingue después de un choque, si y solo si
$R=\beta\rho(W)/\alpha$ es menor que 1.

$R$ funciona como el número reproductivo de una epidemia: compara cuánto déficit genera
la red con cuánto corrige cada entidad. Si $R<1$, la red amortigua. Si $R>1$, la red
amplifica cualquier choque.

La figura muestra la tasa de crecimiento del déficit en el plano $\alpha$–$\beta$ para
las tres topologías. Azul es estable, rojo inestable. La línea negra es la frontera
teórica $\beta=\alpha/\rho(W)$ y coincide con el cambio de color. Las redes más
concentradas tienen $\rho$ mayor y pierden la estabilidad con menos contagio.

Además, los vectores propios dominantes ordenan a las entidades: el izquierdo mide
cuánto déficit propaga cada una (importancia) y el derecho cuánto acumula
(vulnerabilidad).
\end{decir}
\begin{terminos}
\begin{description}
\item[Erdős–Rényi (ER).] Red aleatoria en la que cada arista existe con probabilidad
$p$, independiente de las demás. El número de conexiones por nodo es similar para
todos.
\item[Barabási–Albert (BA).] Red que crece por conexión preferencial: un nodo nuevo se
une a los existentes con probabilidad proporcional a su número de conexiones. Produce
nodos concentradores (\emph{hubs}).
\item[Núcleo-periferia.] Red con un núcleo denso y una periferia que se conecta casi
solo al núcleo.
\item[Distribución lognormal.] Distribución de una variable cuyo logaritmo es normal.
Montos $W_{ij}=\sqrt{s_is_j}\,\varepsilon_{ij}$, con tamaños $s$ y ruido $\varepsilon$
lognormales; luego se normaliza a peso medio 1.
\item[$I$.] Matriz identidad.
\item[Valor propio y vector propio.] $\lambda$ y $v\neq0$ tales que $Wv=\lambda v$. Los
valores propios de $-\alpha I+\beta W$ son $-\alpha+\beta\lambda_k(W)$.
\item[Teorema de Perron–Frobenius.] Una matriz con entradas no negativas tiene un valor
propio real igual a su radio espectral, con un vector propio de entradas no negativas.
\item[$\rho(W)$.] Radio espectral: el mayor módulo de los valores propios de $W$.
\item[$R$.] Índice de amplificación $\beta\rho(W)/\alpha$. Adimensional.
\item[Número reproductivo ($R_0$).] En epidemiología, número medio de contagios que
produce un infectado. La analogía es exacta en forma: umbral en 1.
\item[Importancia y vulnerabilidad.] Componentes de los vectores propios izquierdo
($u^\top W=\rho u^\top$) y derecho ($Wv=\rho v$).
\end{description}
\end{terminos}
\begin{preguntas}
\pr{¿Cómo se verificó?}{Se comparó la integración numérica con la solución exacta
$x(t)=e^{Mt}x_0$; el error relativo es menor que $10^{-6}$. Con $R=0.9$ el déficit se
extingue y con $R=1.1$ crece, como predice la teoría.}
\end{preguntas}

\diapo{8}{Tarea 2: redistribución y estabilidad}{1.5 min}
\begin{decir}
La segunda tarea agrega la redistribución de liquidez. Usamos el Laplaciano de flujo
$L=D_\text{out}-W$, donde $D_\text{out}$ es una matriz diagonal con lo que paga cada
entidad. El término $-D_0Lx$ traslada déficit de una entidad a otra a lo largo de los
pagos.

Las columnas de $L$ suman cero. Eso implica que la suma de todos los déficits no cambia
por este término: el déficit se mueve, pero no se crea ni se destruye. Lo verificamos
numéricamente: la suma cambia menos de $10^{-10}$.

Con los tres términos, la matriz del sistema es $M=-\alpha I-D_0L+\beta W$. Su valor
propio con mayor parte real, $\sigma$, decide la estabilidad: si $\sigma<0$ la red se
recupera, y el tiempo de recuperación es $\tau_\text{rec}=-1/\sigma$.

La figura de la izquierda muestra que las redes con mayor conectividad algebraica
$\lambda_2$ se recuperan más rápido. En redes simétricas la tasa de recuperación es
exactamente $\alpha+D_0\lambda_2$, y la simulación coincide con un error del 1~\%.

La figura de la derecha muestra comunidades: grupos de entidades que se pagan mucho
entre sí y poco con el resto. Las obtenemos con el vector de Fiedler del Laplaciano
normalizado.
\end{decir}
\begin{terminos}
\begin{description}
\item[$D_\text{out}$.] Matriz diagonal con $\mathrm{out}_j=\sum_iW_{ij}$, el pago total
de $j$ por ciclo.
\item[Laplaciano de flujo $L$.] $L=D_\text{out}-W$. Como $\sum_iL_{ij}=\mathrm{out}_j-
\sum_iW_{ij}=0$, se cumple $\mathbf 1^\top L=0$ y $\frac{d}{dt}\sum_ix_i=0$ bajo este
término.
\item[$D_0$.] Velocidad de la redistribución.
\item[$M$.] Matriz del sistema lineal completo.
\item[$\sigma$.] Máxima parte real de los valores propios de $M$. Negativa: estable.
\item[$\tau_\text{rec}$.] Tiempo de recuperación $-1/\sigma$, en ciclos: tiempo en que
una perturbación se reduce en un factor $e$.
\item[$\lambda_2$ (conectividad algebraica).] Segundo valor propio más pequeño del
Laplaciano de la red simetrizada $S=(W+W^\top)/2$. Vale 0 si la red está partida en
dos y crece con la conectividad.
\item[Vector de Fiedler.] Vector propio asociado a $\lambda_2$. El signo de sus
componentes divide la red en dos grupos poco conectados.
\item[Laplaciano normalizado.] $I-D^{-1/2}SD^{-1/2}$, con $D$ la matriz de grados de
$S$. Implementa el corte normalizado de Shi y Malik (2000), que evita separar solo
nodos periféricos.
\item[Bisección recursiva.] Se divide en dos la comunidad más grande hasta obtener el
número pedido de comunidades.
\end{description}
\end{terminos}
\begin{preguntas}
\pr{¿Qué representa la redistribución en la práctica?}{Mecanismos que trasladan presión
de liquidez sin crearla: líneas de crédito interbancario, compensación de pagos. Es una
idealización.}
\end{preguntas}

\diapo{9}{Tarea 3: sincronización de los ciclos de pago}{1.5 min}
\begin{decir}
La tercera tarea modela el ritmo de los pagos con el modelo de Kuramoto, el modelo
estándar de sincronización en física. Cada entidad es un oscilador con una fase
$\theta_i$ y un ritmo natural $\omega_i$. Si el acoplamiento $K$ es grande, las fases se
alinean, como cuando un público que aplaude termina aplaudiendo al unísono.

Medimos la sincronización con el parámetro de orden $r$, que vale 0 si las fases están
dispersas y 1 si todas coinciden. Para cada entidad definimos la coherencia local
$r_i$: el promedio del coseno del desfase con sus pagadores, ponderado por lo que le paga
cada uno.

Este término se conecta con el déficit mediante $\gamma k_i(1-r_i)$. Si el ingreso
llega en fase, $r_i=1$ y el término vale cero. Si llega en contrafase, $r_i=-1$ y el
término es máximo. Así la sincronización deja de ser descriptiva y afecta la liquidez.

Para validar, usamos un resultado exacto: en una red completa con frecuencias de
Lorentz, $r=\sqrt{1-2\Delta/K}$. Con $K=2$ el valor teórico es 0.7071 y el simulado
0.7076. La figura muestra que las redes dispersas se sincronizan de forma más gradual
que la red completa.

Una decisión técnica: normalizamos el acoplamiento por el ingreso medio y no por el
número de entidades. En redes dispersas, dividir por $N$ anula el acoplamiento cuando la
red crece.
\end{decir}
\begin{terminos}
\begin{description}
\item[Modelo de Kuramoto.] $\dot\theta_i=\omega_i+\frac{K}{\langle k\rangle}\sum_jW_{ij}
\sin(\theta_j-\theta_i)$. Cada oscilador acelera o frena según la diferencia de fase con
quienes lo influyen.
\item[$\omega_i$ y distribución de Lorentz.] Frecuencias naturales sorteadas de una
distribución de Lorentz (o de Cauchy) con centro $\bar\omega=2\pi$ (un ciclo por día) y
semiancho $\Delta$.
\item[$\Delta$.] Dispersión de los ritmos naturales. Mayor $\Delta$, más difícil
sincronizar.
\item[$r$.] Parámetro de orden global $r=\bigl|\frac1N\sum_je^{i\theta_j}\bigr|\in[0,1]$.
Aquí $i$ dentro del exponente es la unidad imaginaria.
\item[$r_i$.] Coherencia local $\frac1{k_i}\sum_jW_{ij}\cos(\theta_j-\theta_i)\in[-1,1]$.
\item[$\gamma k_i(1-r_i)$.] Ingreso desfasado convertido en déficit; es igual a
$\gamma\sum_jW_{ij}(1-\cos(\theta_j-\theta_i))$. Unidades de dinero por ciclo, como
los demás términos.
\item[$K_c=2\Delta$.] Acoplamiento crítico en la red completa con Lorentz: por debajo,
$r\approx0$.
\item[Normalización por $\langle k\rangle$.] Rodrigues et al.\ (2016). En la red
completa coincide con el $K/N$ clásico.
\end{description}
\end{terminos}
\begin{preguntas}
\pr{¿Por qué Kuramoto y no otro modelo?}{Es el modelo mínimo con una transición de
sincronización conocida analíticamente. Permite validar el código contra un resultado
exacto y tiene una interpretación directa: cada entidad ajusta su ciclo al de quienes le
pagan.}
\end{preguntas}

\diapo{10}{Tarea 4: impagos y cascadas}{1.5 min}
\begin{decir}
La cuarta tarea introduce el impago. Cada entidad tiene un colchón proporcional a lo que
paga, con algo de variación aleatoria. Si su déficit supera el colchón, la entidad deja
de pagar. Sus receptores pierden ese ingreso, su déficit salta, y algunos pueden cruzar
su propio umbral. Repetimos el proceso en el mismo instante hasta que no hay impagos
nuevos: eso es una cascada.

Simulamos dos choques. El retiro de liquidez es un salto en el déficit de una entidad.
La caída operativa es una entidad que recibe pero no paga durante un tiempo.

Comparamos con tres modelos clásicos de contagio financiero: Eisenberg–Noe, DebtRank y
Watts.

La figura muestra la ventana de contagio. Con pocas conexiones, el choque no tiene por
dónde propagarse. Con muchas conexiones, cada entidad recibe de muchos pagadores y la
pérdida de uno se diluye. Las cascadas son máximas con conectividad intermedia. Nuestra
dinámica produce esa ventana y difiere de DebtRank y de Watts, que no tienen dinámica
temporal ni recuperación.
\end{decir}
\begin{terminos}
\begin{description}
\item[$x_i^{(c)}=c\,\mathrm{out}_i\,\varepsilon_i$.] Umbral: $c$ es el colchón medido en
ciclos de pagos salientes, $\mathrm{out}_i$ lo que paga $i$ por ciclo y $\varepsilon_i$
un factor lognormal de media 1.
\item[$W_{\cdot i}\leftarrow0$.] La columna $i$ de $W$ se hace cero: $i$ deja de pagar
a todos.
\item[$x_k\mathrel{+}=\lambda\,\mathrm{out}_k W_{ki}/k_k$.] Salto en el déficit de cada
receptor $k$: fracción de su ingreso que perdió ($W_{ki}/k_k$) por su escala de pagos
($\mathrm{out}_k$) por $\lambda$, el parámetro de pérdida.
\item[Cascada instantánea.] Repetición de la regla de impago en el mismo instante hasta
que no hay impagos nuevos.
\item[Orden en cada paso.] Eventos programados, paso de integración, detección de
impagos y cascada.
\item[Eisenberg–Noe (2001).] Modelo de liquidación: cada entidad paga el mínimo entre lo
que debe y lo que tiene, $p=\min(\bar p,e+\Pi^\top p)$. $\bar p$ son las obligaciones,
$e$ los activos externos y $\Pi$ la matriz de proporciones de deuda.
\item[DebtRank.] Battiston et al.\ (2012). Propaga estrés de acreedor a deudor y cada
nodo transmite una sola vez, sin reverberación.
\item[Modelo de Watts (2002).] Un nodo falla cuando la fracción de vecinos caídos supera
un umbral $\phi$.
\item[Colapso global.] Impago de al menos el 30~\% de las entidades.
\end{description}
\end{terminos}

% ===================================================================
\section{Tarea 5: predicción sobre simulaciones}

\diapo{11}{Diseño del experimento numérico}{1 min}
\begin{decir}
Una simulación es precisa pero lenta para responder en tiempo real. La quinta tarea
entrena modelos de aprendizaje automático sobre miles de simulaciones para predecir el
resultado en milisegundos.

Generamos 300 redes y 20\,000 escenarios. Cada escenario combina una red con valores de
13 parámetros, entre ellos los de la tabla, el tamaño y el tipo de choque. Los valores se
eligen con un hipercubo latino, que cubre el espacio de parámetros de manera uniforme.
Además, concentramos más escenarios cerca de $R=1$, donde el comportamiento cambia y
predecir es más difícil.

Construimos dos tablas: una por escenario, con la etiqueta de colapso global, y una por
entidad, con el tamaño de la cascada que produce un choque en esa entidad.

Para evaluar, usamos GroupKFold por red: ninguna red aparece a la vez en entrenamiento y
en prueba. Si no se hiciera así, el modelo podría memorizar la red y el resultado sería
optimista.
\end{decir}
\begin{terminos}
\begin{description}
\item[Hipercubo latino.] Muestreo en el que el rango de cada parámetro se divide en $N$
intervalos y cada intervalo se usa exactamente una vez. Cubre mejor el espacio que un
muestreo aleatorio simple.
\item[Muestreo cerca de $R=1$.] $R=1+0.6\,\mathrm{sign}(u)|u|^2$, con $u$ uniforme en
$[-1,1]$.
\item[Los 13 parámetros.] $\alpha$, $R$, $D_0$, $K$, $\gamma$, $\Delta$, colchón $c$,
pérdida $\lambda$, tamaño del choque, duración, tipo de choque, entidad chocada y red.
\item[GroupKFold.] Validación cruzada en $k$ particiones (aquí 5) en la que todas las
filas de un mismo grupo (una red) quedan en la misma partición.
\item[Validación cruzada.] Se entrena con $k-1$ particiones y se evalúa en la restante,
rotando; cada predicción se hace sobre datos que el modelo no vio.
\item[Semilla por fila.] Cada escenario tiene su propia semilla derivada con
\texttt{SeedSequence}: el conjunto de datos se regenera idéntico con cualquier número de
procesos.
\end{description}
\end{terminos}

\diapo{12}{Tarea 5: resultados}{1.5 min}
\begin{decir}
Probamos cuatro modelos. La regresión logística es la referencia interpretable. XGBoost
combina cientos de árboles de decisión. La GCN es una red neuronal que trabaja sobre el
grafo. La PINN es una red neuronal que respeta las ecuaciones del modelo.

Para el colapso global, XGBoost alcanza un AUC de 0.98. Para el tamaño de la cascada
por entidad, la GCN explica el 60~\% de la variación. La figura muestra el AUC del
colapso por tramos de $R$ y por topología. XGBoost se mantiene entre 0.97 y 0.99 en todos
los grupos. La regresión logística baja a 0.92 en la zona más inestable, $R>1.2$, porque
allí la relación entre las variables y el colapso deja de ser lineal.

Una simulación tarda 0.08 segundos; una predicción, 0.2 milisegundos. Eso permite el
mapa de riesgo instantáneo de la aplicación.

La PINN resuelve el problema inverso: a partir de trayectorias observadas, estima los
parámetros. Recupera $\gamma$ y $K$ con error menor al 1~\%. El parámetro $\beta$ tiene
más error, 23~\%, porque su efecto se confunde con el de $\alpha$.

Una limitación: dentro de una misma red, el orden de las entidades por riesgo se predice
solo de forma moderada, con una correlación de Spearman de 0.3. Por eso la aplicación
permite verificar cualquier predicción con una simulación.
\end{decir}
\begin{terminos}
\begin{description}
\item[Regresión logística.] Modelo $P(y=1)=1/(1+e^{-(b_0+\sum_kb_kz_k)})$. Cada
coeficiente indica cuánto cambia el logaritmo de la razón de probabilidades por unidad
de la variable.
\item[XGBoost.] \emph{Extreme Gradient Boosting}. Suma de árboles de decisión entrenados
en secuencia; cada árbol corrige los errores de los anteriores (\emph{gradient
boosting}).
\item[Árbol de decisión.] Modelo que divide los datos con preguntas del tipo ``¿$R>1$?''
hasta llegar a una predicción.
\item[GCN.] \emph{Graph Convolutional Network} (Kipf y Welling, 2017). Cada capa promedia
las variables de los vecinos de cada nodo y aplica una transformación aprendida. Se
usan 3 capas con una conexión residual que lleva las variables propias del nodo hasta
la salida.
\item[PINN.] \emph{Physics-Informed Neural Network} (Raissi et al., 2019). Red neuronal
entrenada para ajustar los datos y, a la vez, cumplir las ecuaciones del modelo. Los
parámetros físicos son incógnitas del entrenamiento.
\item[Problema inverso.] Estimar parámetros a partir de observaciones de la solución.
\item[AUC.] Área bajo la curva ROC. Probabilidad de que el modelo asigne más riesgo a un
caso positivo elegido al azar que a uno negativo. 0.5 es azar; 1 es perfecto.
\item[Curva ROC.] Tasa de verdaderos positivos contra tasa de falsos positivos para todos
los umbrales de decisión.
\item[PR-AUC.] Área bajo la curva precisión-exhaustividad. Más informativa cuando los
positivos son pocos.
\item[$R^2$.] Coeficiente de determinación: fracción de la varianza de la variable
objetivo que explica el modelo.
\item[Spearman ($\rho_s$).] Correlación entre los órdenes (rangos) de dos variables.
\end{description}
\end{terminos}
\begin{preguntas}
\pr{¿Para qué predecir si se puede simular?}{Para explorar en tiempo real: la predicción
es unas 400 veces más rápida por escenario y unas 20\,000 veces más rápida para la red
completa de entidades. La simulación queda como verificación.}
\pr{¿Por qué $\beta$ tiene más error en la PINN?}{En la parte lineal, $\alpha$ y $\beta$
entran combinadas en $M$: varias parejas producen trayectorias casi iguales. Es un
problema de identificabilidad, no del código.}
\end{preguntas}
